\scanpage{72}
\mspage{68}
% source lines 1--5
et on \deleted{pourra ainsi} l'étend en $\alpha_*$ par
\[
\alpha_*(S)=\underset{x\in S}{\operatorname{Sup}\text{ dans }\Sigma_{E'}}
(\underbrace{\alpha(x)}_{\substack{\Vert\\\alpha_*(S_{\deleted{\unsure{E}}}(x))}})
=\operatorname{supp}^{\circ}_{E'}(\bar\alpha(S))
\]
\uline{formule valable}, puisque
\[
S=\left[\underset{x\in S}{\operatorname{Sup}_{\deleted{\unsure{E}}}(\text{dans }\Sigma_E)}\right](S_x).
\]
Ainsi

% source lines 6--11
\uline{Lemme} 1 Corr. biunivoque entre \deleted{\unsure{les}} \deleted{Appl} $\operatorname{Hom\,ens}(E,\Sigma_{E'})$ et $\operatorname{Hom\,ord}^{!}(\Sigma_E,\Sigma_{E'})$, où le $!$ signifie qu'on prend les hom. d'ens. ordonnés commutant aux $\operatorname{Sup}$ quelconques.
% margin note written slanted
\margin{Mais il n'est pas clair qu'on ait la formule, pour tt $A\subset E$ $\deleted{\unsure{\beta|\Delta}}\alpha_*(\operatorname{supp}^{\circ}_E(A))\overset{?}{=}\operatorname{supp}^{\circ}_{E'}(\bar\alpha(A))$ mais seulement que $\supset$}

% source lines 12--18
On voudrait cependant une notion de \unsure{transposé|transposée} d'un tel \unsure{fourbis}, qui généralise le cas connu des ``disjoncteurs triviaux'', où la définition du transposé $\beta$ de $\alpha$ peut se décrire par
\[
\deleted{\alpha(x)\ne}
\]
\[
x'\unsure{\in|\notin}\alpha(x)\Longleftrightarrow\unsure{\alpha|x}\in\beta(x')
\]
% source lines 19--21
ici on passe à la négation
\[
x'\mid\circ\mid\alpha(x)\Longleftrightarrow x\unsure{\mid\circ\mid}\beta(x')
\]
On a donc
% archive sheet number: 71
